\documentclass[10pt,a4paper]{amsart}
\usepackage[margin=19mm]{geometry}
\usepackage{amsmath,amssymb,mathtools}
\usepackage[scr=rsfs,cal=euler]{mathalfa}
\usepackage{xfrac}
\usepackage[unicode,hidelinks,bookmarks=false]{hyperref}
\newcommand{\ie}{i.e.\ }
\newcommand{\cf}{cf.\ }
\newcommand{\resp}{respectively\ }
\newcommand{\Subsection}{\subsection}
\providecommand{\nobreakdash}{\nobreak\hbox{-}}
\setlength{\emergencystretch}{2em}
\multlinegap0pt
\newtheorem{theorem}{Theorem}
\newcommand{\A}{\mathcal{A}}
\newcommand{\C}{\mathbb{C}}
\newcommand{\D}{\mathbb{D}}
\newcommand{\IntComp}{_{\ast}} % Notation for the interior of the complement
\newcommand{\QD}{\text{QD}}
\newcommand{\dEquals}{\dot{=}}
\title{Analysis of Log-Weighted Quadrature Domains}
\author{Andrew Graven}
\date{Source: arXiv:2604.10394v1 (CC BY 4.0)}
\begin{document}
\maketitle

\noindent\emph{Source and licence.} Analysis of Log-Weighted Quadrature Domains. Version arXiv:2604.10394v1 (CC BY 4.0), distributed by arXiv under the Creative Commons Attribution 4.0 International licence, http://creativecommons.org/licenses/by/4.0/, as recorded in the arXiv licence metadata for this exact version and shown on the abstract page https://arxiv.org/abs/2604.10394v1. Full licence notice: \texttt{licenses/arxiv-2604.10394-CC-BY-4.0.txt}.\par\smallskip

\noindent\textit{Editorial scope.} Counted unit: the article's theorem thm:LQDCEZero (source0.tex lines 301--307). The article's complete original proof of it is delivered with it (source0.tex lines 346--390, one \texttt{proof} environment of 45 lines, which the article places away from the statement). No mathematics is written, repaired or re-derived: the statement and the proof are the article's own lines, in the article's own notation, and no retained line is substituted or re-flowed.  No further source-local context is needed: every cross-reference of every retained line resolves inside the two delivered spans.  Nothing was omitted from any retained span, so the package carries no omission record.  No retained line contains a citation, so the wrapper carries no bibliography and no bibliography entry.  No retained line contains a figure environment, an image-inclusion command or drawing code, so no image is delivered and the package declares no asset dependency.  Editorial formatting only, in addition: an AMS \texttt{amsart} preamble that restates the article's own declarations and macros used by the retained lines, namely the environments \texttt{theorem} on separate counters, together with the standard packages those lines need.  The retained environments print in this document as Theorem 1 in order of appearance, whereas the article prints the same items with the numbering of its own full document.  The complete native source of the article as distributed in the arXiv e-print for this exact version is bundled unchanged as \texttt{sources/198270/source0.tex}.
\begin{theorem}\label{thm:LQDCEZero}
    If $\Omega\in\QD_0(h)$ is a domain containing zero, then there exist unique $q\in\C$ and $G\in\A(\Omega)$ such that
    \begin{equation}\label{eqn:LQDCEZero}
        \dfrac{\ln|w|^2}{w}\dEquals h(w)+\dfrac{q}{w}+G(w),
    \end{equation}
    where $G=C_{\rho_0}^{\Omega\IntComp}$.
\end{theorem}

\begin{proof}[Proof of Theorem \ref{thm:LQDCEZero}]\label{proof:LQDCEZero}
Fix $0\in\Omega\in\QD_0(h)$ and $r<d(0,\partial\Omega)$.
We first consider the case in which $\Omega$ is bounded. If $w\in\Omega\IntComp$ then, applying the quadrature identity to $\left(\xi\mapsto\frac{\xi}{w-\xi}\right)\in L_a^1(\Omega;\rho_0)$ yields
\begin{align*}
    \int_{\Omega}\dfrac{\xi|\xi|^{-2}}{\xi-w}dA(\xi)&=\oint_{\partial\Omega}\dfrac{\xi h(\xi)}{\xi-w}d\xi=-wh(w).
\end{align*}
On the other hand, if $w\in\Omega$ then, by Green's theorem ($\Omega\in\QD_0$ by assumption, so its boundary is rectifiable), 
\begin{align*}
    \int_{\Omega}\dfrac{\xi|\xi|^{-2}}{\xi-w}dA(\xi)&=-\ln|w|^2+\oint_{\partial\Omega}\dfrac{\ln|\xi|^2}{\xi-w}d\xi=-\ln|w|^2+w\oint_{\partial\Omega}\dfrac{\ln|\xi|^2}{\xi(\xi-w)}d\xi+\oint_{\partial\Omega}\dfrac{\ln|\xi|^2}{\xi}d\xi.
\end{align*}
Applying Green's theorem to the first integral, and splitting the second to isolate the pole at $0$, we obtain
\begin{align*}
    \int_{\Omega}\dfrac{\xi|\xi|^{-2}}{\xi-w}dA(\xi)&=-\ln|w|^2-w\int_{\Omega\IntComp}\dfrac{|\xi|^{-2}}{\xi-w}d\xi+\left(\oint_{\partial(\Omega\setminus\D_r)}\dfrac{\ln|\xi|^2}{\xi}d\xi+\oint_{\partial\D_r}\dfrac{\ln|\xi|^2}{\xi}d\xi\right)
\end{align*}
We recognize the first integral as a weighted Cauchy transform, and note that the second two integrals integrate to a constant depending only on $\Omega$:
\begin{align*}
    \int_{\Omega}\dfrac{\xi|\xi|^{-2}}{\xi-w}dA(\xi)&=-\ln|w|^2+wC_{\rho_0}^{\Omega\IntComp}(w)+q_{\Omega}.
\end{align*}
where $q_\Omega=\int_{\Omega\setminus\D_r}|\xi|^{-2}dA(\xi)+\ln|r|^2$. Hence,
\begin{align*}
    \dfrac{\ln|w|^2}{w}&\dEquals h(w)+C_{\rho_0}^{\Omega\IntComp}(w)+\dfrac{q_{\Omega}}{w}
\end{align*}

Finally, we consider the case in which $\Omega$ is unbounded. Fix $w_0\in\Omega\IntComp$. If $w\in\Omega\IntComp$ then, applying the quadrature identity to $\frac{\xi}{(w-w_0)(w-\xi)}\in L_a^1(\Omega;\rho_0)$ yields
\begin{align*}
    \int_{\Omega}\dfrac{\xi|\xi|^{-2}}{(\xi-w_0)(\xi-w)}dA(\xi)&=\oint_{\partial\Omega}\dfrac{\xi h(\xi)}{(\xi-w_0)(\xi-w)}d\xi=-\dfrac{wh(w)-w_0h(w_0)}{w-w_0}.
\end{align*}
On the other hand, if $w\in\Omega$ then, by Green's theorem,
\begin{align*}
    \int_{\Omega}\dfrac{\xi|\xi|^{-2}}{(\xi-w_0)(\xi-w)}dA(\xi)&=-\dfrac{\ln|w|^2}{w-w_0}+\oint_{\partial\Omega}\dfrac{\ln|\xi|^2}{(\xi-w_0)(\xi-w)}d\xi\\
    &=\dfrac{-1}{w-w_0}\left(\ln|w|^2-w\oint_{\partial\Omega}\dfrac{\ln|\xi|^2}{\xi(\xi-w)}d\xi+w_0\oint_{\partial\Omega}\dfrac{\ln|\xi|^2}{\xi(\xi-w_0)}d\xi\right)
\end{align*}
Applying Green's theorem to the first integral, and splitting the second to isolate the pole at $0$, we obtain
\begin{align*}
    \int_{\Omega}\dfrac{\xi|\xi|^{-2}}{(\xi-w_0)(\xi-w)}dA(\xi)&=\dfrac{-1}{w-w_0}\left(\ln|w|^2+w\int_{\Omega\IntComp}\dfrac{|\xi|^{-2}}{\xi-w}dA(\xi)+w_0\left(\oint_{\partial(\Omega\setminus\D_r)}\dfrac{\ln|\xi|^2d\xi}{\xi(\xi-w_0)}+\oint_{\partial\D_r}\dfrac{\ln|r|^2d\xi}{\xi(\xi-w_0)}\right)\right)
\end{align*}
We recognize the first integral as a weighted Cauchy transform, the second integral is also recognized as a weighted Cauchy transform after applying Green's theorem, and the third integral may be calculated directly via the residue theorem.
\begin{align*}
    \int_{\Omega}\dfrac{\xi|\xi|^{-2}}{(\xi-w_0)(\xi-w)}dA(\xi)&=\dfrac{-1}{w-w_0}\left(\ln|w|^2-wC_{\rho_0}^{\Omega\IntComp}(w)+w_0\left(\int_{\Omega\setminus\D_r}\dfrac{|\xi|^{-2}}{\xi-w_0}dA(\xi)-\dfrac{\ln|r|^2}{w_0}\right)\right)\\
    &=\dfrac{-1}{w-w_0}\left(\ln|w|^2-wC_{\rho_0}^{\Omega\IntComp}(w)-w_0C_{\rho_0}^{\Omega\setminus\D_r}(w_0)+\ln|r|^2\right).
\end{align*}
Hence for each $w\in\partial\Omega$,
$$\dfrac{\ln|w|^2}{w}\dEquals h(w)+C_{\rho_0}^{\Omega\IntComp}(w)+\dfrac{q_{\Omega,w_0}}{w},$$
where $q_{\Omega,w_0}=w_0C_{\rho_0}^{\Omega\setminus\D_r}(w_0)-w_0h(w_0)+\ln|r|^2$. For the uniqueness of $G$ and $q$, suppose there existed two pairs $G_1,q_1$ and $G_2,q_2$ for which Equation \ref{eqn:LQDCEZero} holds. Subtracting these from each other, and multiplying through by $w$, we find that $wG_1(w)+q_1\dEquals wG_2(w)+q_2$. Both sides of the equation are analytic in $\Omega$ and extend continuously to the boundary, it follows from the identity theorem that $wG_1(w)+q_1=wG_2(w)+q_2$ on $\Omega$. Plugging in $w=0$, we find that $q_1=q_2$, from which it follows that $G_1=G_2$ on $\Omega$.
\end{proof}

\end{document}
